\section{Problema de Pesquisa}
\label{problema}

O Problema de Pesquisa define o coração metodológico do trabalho. Trata-se da delimitação precisa de um desafio técnico, prático ou teórico cuja solução ainda não está satisfatoriamente documentada na literatura científica ou nas práticas consolidadas de mercado.

Uma boa formulação de problema deve atender aos seguintes requisitos científicos:
\begin{enumerate}
    \item \textbf{Clareza e Precisão:} O problema deve ser exposto de forma objetiva, evitando ambiguidades ou termos vagos.
    \item \textbf{Delimitação de Escopo:} Deve ser suficientemente circunscrito para permitir uma investigação aprofundada dentro das restrições de tempo e recursos do projeto.
    \item \textbf{Viabilidade de Verificação:} O problema deve ser passível de resposta empírica ou analítica por meio de métodos científicos replicáveis.
\end{enumerate}

Para articular o problema, recomenda-se explicitar a discrepância entre o \textit{estado atual} (o cenário observado com suas ineficiências ou limitações) e o \textit{estado desejado} (o cenário otimizado proposto pela pesquisa). 

Estruturas fraseológicas como \enquote{Apesar dos recentes avanços no domínio X, observa-se que as abordagens existentes ainda enfrentam limitações quanto a Y...} auxiliam na evidenciação da lacuna de conhecimento que a dissertação ou tese se propõe a superar.